% =============================================================================
% Modern C++ Masterclass — Power Systems Edition
% Beamer Presentation for Universitas Bengkulu
% Compile with: latexmk -pdf -shell-escape modern_cpp_slides.tex
% =============================================================================
\documentclass[aspectratio=169,10pt]{beamer}

% Theme
\usepackage{beamerthemeUB}

% Packages
\usepackage{amsmath,amssymb,amsthm}
\usepackage{mathtools}
\usepackage{booktabs}
\usepackage{graphicx}
\usepackage{hyperref}
\usepackage{listings}
\usepackage{xcolor}
\usepackage{multicol}
\usepackage{textcomp}

% Snippet commands - direct input (no index file to avoid preamble issues)
% Usage: \begin{minted}[linenos,breaklines,frame=lines,framesep=2mm]{cpp}
consteval int compile_time_fibonacci(int n) {
    if (n <= 0) return 0;
    if (n == 1) return 1;
    int a = 0;
    int b = 1;
    for (int i = 2; i <= n; ++i) {
        int next = a + b;
        a = b;
        b = next;
    }
    return b;
}
\end{minted}


% Custom commands
\newcommand{\cpp}{\texttt{C++}}
\newcommand{\cpptwenty}{\texttt{C++20}}
\newcommand{\cpptwentythree}{\texttt{C++23}}
\newcommand{\consteval}{\texttt{consteval}}
\newcommand{\constexpr}{\texttt{constexpr}}
\newcommand{\stringview}{\texttt{std::string\_view}}
\newcommand{\stdspan}{\texttt{std::span}}
\newcommand{\optional}{\texttt{std::optional}}
\newcommand{\expected}{\texttt{std::expected}}
\newcommand{\jthread}{\texttt{std::jthread}}
\newcommand{\atomic}{\texttt{std::atomic}}
\newcommand{\mutex}{\texttt{std::mutex}}
\newcommand{\scopedlock}{\texttt{std::scoped\_lock}}
\newcommand{\ranges}{\texttt{std::ranges}}
\newcommand{\views}{\texttt{std::views}}
\newcommand{\format}{\texttt{std::format}}
\newcommand{\concepts}{\texttt{concepts}}
\newcommand{\AoS}{\texttt{AoS}}
\newcommand{\SoA}{\texttt{SoA}}
\newcommand{\RAII}{\texttt{RAII}}
\newcommand{\modules}{\texttt{modules}}
\newcommand{\coroutines}{\texttt{coroutines}}
\newcommand{\stoptoken}{\texttt{stop\_token}}

% Title info
\title[Modern C++ Masterclass]{Modern C++ Masterclass: \cpptwenty/\cpptwentythree for Power Systems}
\subtitle{N-1 Contingency Analysis with Modern C++}
\author[Novalio Daratha]{Novalio Daratha}
\institute[UNIB]{Fakultas Teknik, Universitas Bengkulu}
\date[\today]{\today}

\begin{document}

% =============================================================================
% TITLE PAGE
% =============================================================================
\begin{frame}
  \titlepage
\end{frame}

% =============================================================================
% TABLE OF CONTENTS
% =============================================================================
\begin{frame}{Daftar Isi}
  \tableofcontents
\end{frame}

% =============================================================================
% SECTION 1: PENDAHULUAN
% =============================================================================
\section{Pendahuluan}

\begin{frame}{Mengapa Modern C++ untuk Power Systems?}
  \begin{block}{Tantangan Domain Kelistrikan}
    \begin{itemize}
      \item \textbf{Data besar}: Ribuan bus, cabang, pengukuran PMU real-time
      \item \textbf{Performa kritis}: N-1 contingency screening butuh $< 1$ detik
      \item \textbf{Keandalan}: Kecil error numerik $\rightarrow$ blackout
      \item \textbf{Kompleksitas}: Model jaringan, proteksi, optimalisasi daya
    \end{itemize}
  \end{block}
  
  \begin{block}{Solusi Modern C++}
    \begin{itemize}
      \item \cpptwenty/\cpptwentythree: \consteval, \ranges, \concepts, \coroutines, \modules
      \item \textbf{Zero-cost abstractions}: Tidak ada overhead runtime
      \item \textbf{Memory safety}: \RAII, \stdspan, \optional, \expected
      \item \textbf{Concurrency}: \jthread, \atomic, \stoptoken
      \item \textbf{Hardware sympathy}: Cache-friendly \SoA, SIMD-ready
    \end{itemize}
  \end{block}
\end{frame}

\begin{frame}{Struktur Kursus (7 Fase + Capstone)}
  \begin{center}
    \begin{tabular}{cll}
      \toprule
      \textbf{Fase} & \textbf{Topik} & \textbf{Proyek} \\
      \midrule
      1 & \consteval, \stringview, \format, Structured Binding, \modules & Log Parser \\
      2 & \RAII, Move Semantics, \stdspan, Testing (Catch2) & MatrixBuffer \\
      3 & \concepts, Type Erasure, \texttt{if constexpr} & Event Dispatcher \\
      4 & \ranges, \views, Pipelines, \optional & PMU Data Pipeline \\
      5 & \jthread, \stoptoken, \atomic, \scopedlock & Parallel Worker Pool \\
      6 & \expected, \SoA vs \AoS, Benchmark & Hardware Sympathy \\
      Capstone & N-1 Parallel Contingency Engine & \textbf{Final Project} \\
      \bottomrule
    \end{tabular}
  \end{center}
  
  \vspace{0.5cm}
  \begin{alertblock}{Metode: Active Learning}
    Setiap fase: Konsep $\rightarrow$ Demo $\rightarrow$ Latihan (C1--C6) $\rightarrow$ Refleksi
  \end{alertblock}
\end{frame}

% =============================================================================
% SECTION 2: PHASE 1
% =============================================================================
\section{Fase 1: Fondasi Modern C++}

\begin{frame}{Fase 1: \consteval, \stringview, \format, Structured Binding}
  \begin{columns}[T]
    \column{0.55\textwidth}
    \begin{block}{Kunci Pembelajaran}
      \begin{itemize}
        \item \consteval: Jaminan compile-time evaluation
        \item \stringview: Non-owning string slice (zero-copy)
        \item \format: Type-safe formatting (seperti Python f-string)
        \item Structured binding: Unpack struct/tuple clean
        \item \modules: Kompilasi cepat, enkapsulasi header
      \end{itemize}
    \end{block}
    
    \column{0.45\textwidth}
    \begin{exampleblock}{Contoh Domain}
      \begin{itemize}
        \item Parse log grid: \texttt{2026-10-01T08:07:45 [CRITICAL] ...}
        \item Koordinat substation: \texttt{lat, lon}
        \item Pembacaan tegangan: \texttt{voltage\_pu}
      \end{itemize}
    \end{exampleblock}
  \end{columns}
\end{frame}

\begin{frame}[fragile]{Consteval: Fibonacci Compile-Time}
  \begin{minted}[linenos,breaklines,frame=lines,framesep=2mm]{cpp}
consteval int compile_time_fibonacci(int n) {
    if (n <= 0) return 0;
    if (n == 1) return 1;
    int a = 0;
    int b = 1;
    for (int i = 2; i <= n; ++i) {
        int next = a + b;
        a = b;
        b = next;
    }
    return b;
}
\end{minted}

\end{frame}

\begin{frame}[fragile]{String View: Zero-Copy Parsing}
  \begin{minted}[linenos,breaklines,frame=lines,framesep=2mm]{cpp}
void print_slice_info(std::string_view text) {
    std::cout << std::format("Slice length: {}, prefix(4): \"{}\"\n", 
                             text.size(), 
                             text.substr(0, std::min<size_t>(4, text.size())));
}
\end{minted}

  \vspace{0.3cm}
  \begin{minted}[linenos,breaklines,frame=lines,framesep=2mm]{cpp}
    std::string full_log = "2026-10-01T11:45:00 [CRITICAL] Transformer Overload Detected";
    // Passing substring view without copying or allocating new string
    std::string_view log_view = full_log;
    print_slice_info(log_view);
\end{minted}

\end{frame}

\begin{frame}[fragile]{Structured Binding: Unpack Struct}
  \begin{minted}[linenos,breaklines,frame=lines,framesep=2mm]{cpp}
    Measurement m{"Bus-01_Sensor", 230.5, 12.8};
    auto [name, v, i] = m; // unpacks struct fields cleanly
    double power = v * i;
    std::cout << std::format("[2] Structured Binding:\n"
                             "    Sensor: {}\n"
                             "    Voltage: {:.1f} V, Current: {:.1f} A -> Apparent Power: {:.2f} VA\n",
                             name, v, i, power);

    // Structured binding with std::tuple / std::pair
    auto get_coordinates = []() -> std::pair<double, double> {
        return {-0.7893, 100.3542}; // e.g., Padang / Sumatra latitude & longitude
    };
    auto [lat, lon] = get_coordinates();
    std::cout << std::format("    Coordinates: ({:.4f}, {:.4f})\n", lat, lon);
\end{minted}

\end{frame}

\begin{frame}[fragile]{Range-for dengan Init Statement}
  \begin{minted}[linenos,breaklines,frame=lines,framesep=2mm]{cpp}
    for (std::vector<double> readings{1.02, 0.99, 1.05, 0.96}; 
         const auto& val : readings) {
        std::cout << std::format("    Voltage pu: {:.2f} ({})\n", 
                                 val, (val >= 0.95 && val <= 1.05) ? "NORMAL" : "VIOLATION");
    }
\end{minted}

\end{frame}

\begin{frame}{Fase 1: C++20 Modules}
  \begin{columns}[T]
    \column{0.5\textwidth}
    \begin{block}{grid\_parser.cppm}
      \begin{itemize}
        \item \texttt{export module grid\_parser;}
        \item \texttt{export enum class Severity}
        \item \texttt{export struct LogEntry}
        \item \texttt{export LogEntry parse\_log\_line(...)}
      \end{itemize}
    \end{block}
    \column{0.5\textwidth}
    \begin{block}{phase1\_modules.cpp}
      \begin{itemize}
        \item \texttt{import grid\_parser;}
        \item Tidak perlu \texttt{\#include}
        \item Kompilasi terpisah, cepat
        \item ODR-safe
      \end{itemize}
    \end{block}
  \end{columns}
  
  \vspace{0.5cm}
  \begin{alertblock}{Catatan: Butuh Clang $\ge$ 17 + CMake $\ge$ 3.28}
  \end{alertblock}
\end{frame}

% =============================================================================
% SECTION 3: PHASE 2
% =============================================================================
\section{Fase 2: Memory Management \& Testing}

\begin{frame}{Fase 2: RAII, Move Semantics, \stdspan, Testing}
  \begin{columns}[T]
    \column{0.55\textwidth}
    \begin{block}{Kunci Pembelajaran}
      \begin{itemize}
        \item \RAII: Resource Acquisition Is Initialization
        \item \texttt{unique\_ptr}: Exclusive ownership
        \item Move semantics: Transfer ownership tanpa copy
        \item \stdspan: Non-owning view, bounds-safe
        \item Catch2: Modern testing framework
      \end{itemize}
    \end{block}
    \column{0.45\textwidth}
    \begin{exampleblock}{MatrixBuffer}
      \begin{itemize}
        \item 2D matrix buffer bounds-checked
        \item \texttt{get\_row()} $\rightarrow$ \stdspan
        \item Non-copyable, movable
        \item Rule of Five + Rule of Zero
      \end{itemize}
    \end{exampleblock}
  \end{columns}
\end{frame}

\begin{frame}[fragile]{RAII dengan unique\_ptr}
  \begin{minted}[linenos,breaklines,frame=lines,framesep=2mm]{cpp}
    std::unique_ptr<T[]> data_{nullptr};
\end{minted}

\end{frame}

\begin{frame}[fragile]{Move Semantics: Transfer Ownership}
  \begin{minted}[linenos,breaklines,frame=lines,framesep=2mm]{cpp}
    MatrixBuffer(MatrixBuffer&& other) noexcept 
        : rows_(std::exchange(other.rows_, 0)),
          cols_(std::exchange(other.cols_, 0)),
          data_(std::move(other.data_)) {}

    MatrixBuffer& operator=(MatrixBuffer&& other) noexcept {
        if (this != &other) {
            rows_ = std::exchange(other.rows_, 0);
            cols_ = std::exchange(other.cols_, 0);
            data_ = std::move(other.data_); // Safely overwrites and frees current memory
        }
        return *this;
    }
\end{minted}

\end{frame}

\begin{frame}[fragile]{\stdspan: Bounds-Safe Row Access}
  \begin{minted}[linenos,breaklines,frame=lines,framesep=2mm]{cpp}
    std::span<T> get_row(size_t row) {
        if (row >= rows_) throw std::out_of_range("Row index out of bounds");
        return std::span<T>(data_.get() + (row * cols_), cols_);
    }
    
    std::span<const T> get_row(size_t row) const {
        if (row >= rows_) throw std::out_of_range("Row index out of bounds");
        return std::span<const T>(data_.get() + (row * cols_), cols_);
    }
\end{minted}

\end{frame}

\begin{frame}[fragile]{Testing dengan Catch2}
  \begin{block}{Test Coverage MatrixBuffer}
    \begin{itemize}
      \item RAII \& Initialization
      \item Move Semantics (construct + assign)
      \item \stdspan prevents buffer decay
    \end{itemize}
  \end{block}
  
  \begin{exampleblock}{Menjalankan Test}
    \begin{minted}{bash}
cmake --preset=release
ninja -C build
ctest --test-dir build --output-on-failure
    \end{minted}
  \end{exampleblock}
  
  \vspace{0.3cm}
  \begin{alertblock}{42 Test Total (Semua Fase)}
    \texttt{test\_matrix}, \texttt{test\_event\_dispatcher}, \texttt{test\_log\_parser}, 
    \texttt{test\_grid\_parser}, \texttt{test\_concurrency}, \texttt{test\_phase6}
  \end{alertblock}
\end{frame}

% =============================================================================
% SECTION 4: PHASE 3
% =============================================================================
\section{Fase 3: Generic Programming \& Concepts}

\begin{frame}{Fase 3: Concepts, Type Erasure, if constexpr}
  \begin{columns}[T]
    \column{0.55\textwidth}
    \begin{block}{Kunci Pembelajaran}
      \begin{itemize}
        \item \concepts: Compile-time contracts
        \item \texttt{requires} clause: Constrain template
        \item Type Erasure: Heterogeneous storage
        \item \texttt{if constexpr}: Compile-time branching
        \item \texttt{std::invocable}: Callable constraint
      \end{itemize}
    \end{block}
    \column{0.45\textwidth}
    \begin{exampleblock}{EventDispatcher}
      \begin{itemize}
        \item \texttt{EventBus} generic
        \item \texttt{LoggableEvent} concept
        \item Auto-logging via \texttt{if constexpr}
        \item Type-safe subscribe/publish
      \end{itemize}
    \end{exampleblock}
  \end{columns}
\end{frame}

\begin{frame}[fragile]{LoggableEvent Concept}
  \begin{minted}[linenos,breaklines,frame=lines,framesep=2mm]{cpp}
    std::string message() const {
        return std::format("CRITICAL: {} has tripped. N-1 constraint violated.", line_id);
    }
\end{minted}

\end{frame}

\begin{frame}[fragile]{Subscribe \& Publish Pattern}
  \begin{minted}[linenos,breaklines,frame=lines,framesep=2mm]{cpp}
    // Subscribe to Voltage drops
    bus.subscribe<VoltageDropEvent>([](const VoltageDropEvent& e) {
        std::cout << std::format("  -> [Controller A] Detected {} voltage drop to {:.2f} pu\n", 
                                 e.substation, e.voltage_pu);
    });

    // Subscribe to Line trips
    bus.subscribe<LineTripEvent>([](const LineTripEvent& e) {
        std::cout << std::format("  -> [Controller B] Dispatching repair crew to {}\n", e.line_id);
    });

    // --- Publishing Events ---
    
    std::cout << "--- Publishing VoltageDropEvent ---\n";
    // This is NOT a LoggableEvent, so 'if constexpr' ignores the auto-log.
    bus.publish(VoltageDropEvent{"Substation_Alpha", 0.92});

    std::cout << "\n--- Publishing LineTripEvent ---\n";
    // This IS a LoggableEvent. The bus will auto-log it before triggering the handlers!
    bus.publish(LineTripEvent{"Line_L42"});
\end{minted}

\end{frame}

\begin{frame}[fragile]{Contoh Output EventDispatcher}
  \begin{minted}{text}
=== Phase 3: Generic Event Dispatcher (C++20 Concepts) ===

--- Publishing VoltageDropEvent ---
  -> [Controller A] Detected Substation_Alpha voltage drop to 0.92 pu

--- Publishing LineTripEvent ---
[SYSTEM LOG] Auto-logged: CRITICAL: Line_L42 has tripped. N-1 constraint violated.
  -> [Controller B] Dispatching repair crew to Line_L42
  \end{minted}
  
  \vspace{0.3cm}
  \begin{itemize}
    \item \texttt{VoltageDropEvent} tanpa \texttt{message()} $\rightarrow$ tidak auto-log
    \item \texttt{LineTripEvent} dengan \texttt{message()} $\rightarrow$ auto-log aktif
    \item Compile-time decision via \texttt{if constexpr}
  \end{itemize}
\end{frame}

% =============================================================================
% SECTION 5: PHASE 4
% =============================================================================
\section{Fase 4: Ranges \& Functional Pipelines}

\begin{frame}{Fase 4: Ranges, Views, Pipelines, Optional}
  \begin{columns}[T]
    \column{0.55\textwidth}
    \begin{block}{Kunci Pembelajaran}
      \begin{itemize}
        \item \ranges + \views: Lazy, composable pipelines
        \item \texttt{filter}, \texttt{transform}, \texttt{take}
        \item Zero intermediate allocation
        \item \optional: Explicit nullable return
        \item \texttt{if (auto x = f(); x.has\_value())}
      \end{itemize}
    \end{block}
    \column{0.45\textwidth}
    \begin{exampleblock}{PMU Data Pipeline}
      \begin{itemize}
        \item Filter invalid readings
        \item Filter voltage $< 0.95$ pu
        \item Transform ke string alert
        \item Take first 2 critical drops
      \end{itemize}
    \end{exampleblock}
  \end{columns}
\end{frame}

\begin{frame}[fragile]{Ranges Pipeline: Lazy Evaluation}
  \begin{minted}[linenos,breaklines,frame=lines,framesep=2mm]{cpp}
inline std::vector<std::string>
critical_drop_alerts(const std::vector<SensorReading>& readings,
                     std::size_t max_alerts = 2) {
    auto alerts = readings
        | std::views::filter([](const SensorReading& r) { return r.is_valid; })  // Step 1: Ignore glitches
        | std::views::filter([](const SensorReading& r) { return r.voltage_pu < 0.95; })  // Step 2: Low voltage
        | std::views::transform([](const SensorReading& r) {  // Step 3: Map to a string
              return std::format("Time {}ms: CRITICAL VOLTAGE {:.2f} pu",
                                 r.time_ms, r.voltage_pu);
          })
        | std::views::take(max_alerts);  // Step 4: Stop after max_alerts

    // Consume the lazy view once into a concrete vector.
    // (take_view's counted_iterator cannot feed std::vector's
    //  iterator-pair constructor directly in C++20.)
    std::vector<std::string> result;
    for (const auto& alert : alerts) {
        result.push_back(alert);
    }
    return result;
}
\end{minted}

\end{frame}

\begin{frame}[fragile]{Optional: Modern Error Handling}
  \begin{minted}[linenos,breaklines,frame=lines,framesep=2mm]{cpp}
inline std::optional<int> find_first_drop(const std::vector<SensorReading>& data) {
    for (const auto& r : data) {
        if (r.is_valid && r.voltage_pu < 0.90) {
            return r.time_ms; // Found it!
        }
    }
    return std::nullopt; // No data matched
}
\end{minted}

\end{frame}

\begin{frame}[fragile]{Contoh Output Pipeline}
  \begin{minted}{text}
--- 1. C++20 Lazy Ranges Pipeline ---
  -> Time 400ms: CRITICAL VOLTAGE 0.92 pu
  -> Time 600ms: CRITICAL VOLTAGE 0.88 pu

--- 2. Modern Error Handling (std::optional) ---
  -> Extreme drop (<0.90) detected at 600ms
  \end{minted}
  
  \vspace{0.3cm}
  \begin{itemize}
    \item Pipeline \textbf{lazy}: Hanya memproses 2 elemen (take(2))
    \item \texttt{transform} hanya dijalankan 2x, bukan 8x
    \item \texttt{optional} memaksa caller handle kasus "not found"
  \end{itemize}
\end{frame}

% =============================================================================
% SECTION 6: PHASE 5
% =============================================================================
\section{Fase 5: Concurrency}

\begin{frame}{Fase 5: jthread, stop\_token, atomic, scoped\_lock}
  \begin{columns}[T]
    \column{0.55\textwidth}
    \begin{block}{Kunci Pembelajaran}
      \begin{itemize}
        \item \jthread: Auto-join pada destructor
        \item \texttt{stop\_token}: Kooperatif cancellation
        \item \atomic: Lock-free counter
        \item \scopedlock: RAII mutex, deadlock-safe
        \item \texttt{memory\_order\_relaxed}: Untuk counter
      \end{itemize}
    \end{block}
    \column{0.45\textwidth}
    \begin{exampleblock}{Parallel Worker Pool}
      \begin{itemize}
        \item 4 worker nodes
        \item Simulasi N-1 contingency calc
        \item Graceful shutdown via \texttt{request\_stop()}
        \item Thread-safe logging + counter
      \end{itemize}
    \end{exampleblock}
  \end{columns}
\end{frame}

\begin{frame}[fragile]{Mutex \& Scoped Lock}
  \begin{minted}[linenos,breaklines,frame=lines,framesep=2mm]{cpp}
std::mutex log_mutex;

void safe_print(const std::string& msg) {
    // std::scoped_lock (C++17) is RAII for locks. 
    // It locks the mutex here, and AUTOMATICALLY unlocks it when the function ends.
    // If an exception occurs, it still unlocks safely. No deadlocks!
    std::scoped_lock lock(log_mutex);
    std::cout << msg << "\n";
}
\end{minted}

\end{frame}

\begin{frame}[fragile]{Atomic Counter (Lock-Free)}
  \begin{minted}[linenos,breaklines,frame=lines,framesep=2mm]{cpp}
std::atomic<int> total_contingencies_analyzed{0};
\end{minted}

\end{frame}

\begin{frame}[fragile]{jthread dengan Stop Token}
  \begin{minted}[linenos,breaklines,frame=lines,framesep=2mm]{cpp}
void worker_node(std::stop_token stoken, int node_id) {
    safe_print(std::format("Node {} online.", node_id));
    
    int local_work = 0;
    
    // The loop checks if the main thread has requested it to stop.
    // This allows the thread to finish its current calculation safely before exiting.
    while (!stoken.stop_requested()) {
        
        // Simulate a complex N-1 contingency calculation (takes 100ms)
        std::this_thread::sleep_for(std::chrono::milliseconds(100));
        
        // Safely increment the global counter without a lock
        total_contingencies_analyzed.fetch_add(1, std::memory_order_relaxed);
        local_work++;
        
        if (local_work % 5 == 0) {
            safe_print(std::format("  -> Node {} completed {} calculations...", node_id, local_work));
        }
    }
    
    safe_print(std::format("Node {} gracefully shutting down.", node_id));
}
\end{minted}

\end{frame}

\begin{frame}[fragile]{jthread RAII: Auto-Join}
  \begin{minted}[linenos,breaklines,frame=lines,framesep=2mm]{cpp}
    {
        // We use std::jthread (Joining Thread) instead of older std::thread
        std::vector<std::jthread> cluster;
        
        // Spawn 4 worker threads
        for (int i = 0; i < 4; ++i) {
            // std::jthread automatically generates a stop_token and passes it to the function!
            cluster.emplace_back(worker_node, i + 1);
        }
        
        safe_print("[Main] All nodes spawned. Letting them run for ~800ms...");
        std::this_thread::sleep_for(std::chrono::milliseconds(800));
        
        safe_print("\n[Main] Requesting all nodes to stop...");
        
        // We gently ask all threads to stop (flips their stop_token to true)
        for (auto& node : cluster) {
            node.request_stop();
        }
        
        safe_print("[Main] End of scope reached. jthreads will automatically join!");
        // --- THE MAGIC OF jthread ---
        // Right here, as 'cluster' goes out of scope, the destructor for every std::jthread 
        // is called. It will AUTOMATICALLY wait (.join()) for the threads to finish cleanly.
        // Old std::thread would just crash the program if you forgot to call .join().
    }
\end{minted}

\end{frame}

\begin{frame}[fragile]{Contoh Output Concurrency}
  \begin{minted}{text}
=== Phase 5: Modern Concurrency (C++20 std::jthread) ===

Node 1 online.
Node 2 online.
Node 3 online.
Node 4 online.
[Main] All nodes spawned. Letting them run for ~800ms...
  -> Node 3 completed 5 calculations...
  -> Node 1 completed 5 calculations...
  -> Node 4 completed 5 calculations...
  -> Node 2 completed 5 calculations...

[Main] Requesting all nodes to stop...
[Main] End of scope reached. jthreads will automatically join!
Node 1 gracefully shutting down.
...

[Main] Final Report: 32 total contingencies analyzed safely.
  \end{minted}
\end{frame}

% =============================================================================
% SECTION 7: PHASE 6
% =============================================================================
\section{Fase 6: Advanced \& Hardware Sympathy}

\begin{frame}{Fase 6: std::expected, SoA vs AoS, Benchmark}
  \begin{columns}[T]
    \column{0.55\textwidth}
    \begin{block}{Kunci Pembelajaran}
      \begin{itemize}
        \item \expected: Error handling tanpa exception
        \item Monadic ops: \texttt{and\_then}, \texttt{transform}, \texttt{or\_else}
        \item \AoS vs \SoA: Memory layout \& cache
        \item Benchmark valid: \texttt{steady\_clock}, \texttt{\textmu s}, anti-DSE
      \end{itemize}
    \end{block}
    \column{0.45\textwidth}
    \begin{exampleblock}{Power Flow Solver}
      \begin{itemize}
        \item \texttt{solve\_power\_flow(load\_mw)} $\rightarrow$ \expected
        \item Error: InvalidTelemetry, ConvergenceFailed, MatrixSingular
        \item Enum error type: Trivial, no heap allocation
      \end{itemize}
    \end{exampleblock}
  \end{columns}
\end{frame}

\begin{frame}[fragile]{std::expected: Error Handling Tanpa Exception}
  \begin{minted}{cpp}
enum class SystemError {
    InvalidTelemetry, ConvergenceFailed, MatrixSingular
};

std::expected<double, SystemError> solve_power_flow(double load_mw) {
    if (load_mw < 0) return std::unexpected(SystemError::InvalidTelemetry);
    if (load_mw > 10000) return std::unexpected(SystemError::ConvergenceFailed);
    return load_mw * 0.95; // Mock calculation
}

// Usage
auto result = solve_power_flow(500.0);
if (result) {
    use(result.value()); // 475.0
} else {
    handle(result.error());
}
  \end{minted}
\end{frame}

\begin{frame}[fragile]{SoA vs AoS: Memory Layout}
  \begin{columns}[T]
    \column{0.5\textwidth}
    \begin{block}{Array of Structs (AoS)}
      \begin{minted}{cpp}
struct BranchOOP {
    int id;
    double r, x, limit;
};
std::vector<BranchOOP> aos_grid;
// Memory: [id][r][x][limit][id][r][x][limit]...
// Stride: 32 bytes, hanya pakai x (8 bytes)
// 75% bandwidth wasted!
      \end{minted}
    \end{block}
    \column{0.5\textwidth}
    \begin{block}{Struct of Arrays (SoA)}
      \begin{minted}{cpp}
struct GridSOA {
    std::vector<int> id;
    std::vector<double> r, x, limit;
};
GridSOA soa_grid;
// Memory: [x][x][x][x]... contiguous!
// Stride: 8 bytes, 100% bandwidth used
// Cache line: 64 bytes = 8 doubles
      \end{minted}
    \end{block}
  \end{columns}
\end{frame}

\begin{frame}{Benchmark: Strided Access (10M Branches)}
  \begin{block}{Hasil pada Clang 23, Apple M-series}
    \begin{center}
      \begin{tabular}{lcc}
        \toprule
        \textbf{Layout} & \textbf{Time} & \textbf{Speedup} \\
        \midrule
        AoS (OOP) & 32,376 $\mu$s & 1.0x (baseline) \\
        SoA (Data-Oriented) & 1,492 $\mu$s & \textbf{21.7x faster} \\
        \bottomrule
      \end{tabular}
    \end{center}
  \end{block}
  
  \begin{alertblock}{Kunci Benchmark Valid}
    \begin{itemize}
      \item \texttt{std::chrono::steady\_clock} + microseconds
      \item Strided access (stride 64) $\rightarrow$ cache miss pada AoS
      \item Anti-DSE: \texttt{asm volatile} consume result
      \item Kondisional claim: hanya claim faster jika speedup $> 1.05$
    \end{itemize}
  \end{alertblock}
\end{frame}

% =============================================================================
% SECTION 8: CAPSTONE
% =============================================================================
\section{Capstone: N-1 Contingency Engine}

\begin{frame}{Capstone: Parallel N-1 Contingency Screening}
  \begin{block}{Problem Statement}
    Diberikan grid $N$ cabang, simulasi \textbf{N-1 contingency} untuk setiap cabang:
    \begin{enumerate}
      \item Trip cabang $i$
      \item Hitung DC Power Flow (Y-bus, line flows)
      \item Cek thermal limit (overload $> 100\%$)
      \item Deteksi islanding (loss of load)
      \item Rank berdasarkan severity
    \end{enumerate}
  \end{block}
  
  \begin{block}{Solusi Modern C++}
    \begin{itemize}
      \item \jthread pool (2 threads, split via \stdspan)
      \item \atomic counter untuk progress tracking
      \item \ranges pipeline untuk filter/sort hasil
      \item \format untuk output terstruktur
    \end{itemize}
  \end{block}
\end{frame}

\begin{frame}[fragile]{Capstone Output}
  \begin{minted}{text}
========================================================
  MODERN C++ CAPSTONE: Parallel N-1 Contingency Engine  
========================================================

[System] Loaded 8 branches into memory.
[System] Spawning thread pool for N-1 analysis...
[System] Analysis complete. 8 contingencies processed in 618 ms.

--- TOP CRITICAL CONTINGENCIES (THERMAL OVERLOAD) ---
 -> Tripping L1 causes cascaded OVERLOAD of 135.5%
 -> Tripping L7 causes cascaded OVERLOAD of 135.5%
 -> Tripping L4 causes cascaded OVERLOAD of 135.5%

--- ISLANDING CONTINGENCIES (LOSS OF LOAD) ---
 -> Tripping L5 isolates buses, causing grid islanding.
 -> Tripping L8 isolates buses, causing grid islanding.
  \end{minted}
\end{frame}

% =============================================================================
% SECTION 9: KESIMPULAN
% =============================================================================
\section{Kesimpulan}

\begin{frame}{Ringkasan: Modern C++ untuk Power Systems}
  \begin{columns}[T]
    \column{0.5\textwidth}
    \begin{block}{Fitur Modern C++ Terpakai}
      \begin{itemize}
        \item \consteval, \concepts, \modules (C++20)
        \item \ranges, \views, \optional, \format (C++20)
        \item \jthread, \stoptoken, \atomic, \scopedlock (C++20)
        \item \expected, \expected monadic ops (C++23)
      \end{itemize}
    \end{block}
    \column{0.5\textwidth}
    \begin{block}{Best Practices}
      \begin{itemize}
        \item RAII everywhere (unique\_ptr, span, jthread)
        \item Value semantics + move, avoid copy
        \item Compile-time checks (concepts, consteval)
        \item Hardware-aware (SoA, cache-friendly)
        \item Test-driven (Catch2, 42 tests pass)
      \end{itemize}
    \end{block}
  \end{columns}
\end{frame}

\begin{frame}{Referensi \& Lanjut}
  \begin{block}{Standar \& Referensi}
    \begin{itemize}
      \item ISO C++20/23 Standard
      \item \textit{C++ Templates} — Vandevoorde/Josuttis
      \item \textit{Effective Modern C++} — Scott Meyers
      \item CppReference.com
    \end{itemize}
  \end{block}
  
  \begin{block}{Topik Lanjut}
    \begin{itemize}
      \item C++26: Reflection, Contracts, std::execution
      \item SYCL/oneAPI untuk GPU acceleration
      \item Coroutines untuk async power flow
      \item Formal verification (CBMC, KLEE)
    \end{itemize}
  \end{block}
\end{frame}

\begin{frame}{Terima Kasih}
  \begin{center}
    \vspace{2cm}
    {\LARGE \textbf{Terima Kasih}}
    
    \vspace{1cm}
    {\large Pertanyaan \& Diskusi}
    
    \vspace{2cm}
    \includegraphics[height=1.5cm]{assets/ub_logo.pdf}
    
    \vspace{0.5cm}
    \texttt{novalio.daratha@unib.ac.id}
  \end{center}
\end{frame}

% =============================================================================
% APPENDIX: PROBLEM SET (C1-C6)
% =============================================================================
\appendix
\section{Lampiran: Problem Set Bloom C1--C6}

\begin{frame}[allowframebreaks]{Problem Set: Fase 1 (C1--C6)}
  \frametitle{Problem Set: Fase 1 — Fondasi Modern C++}
  
  \textbf{C1 (Remember):} Sebutkan 3 perbedaan \texttt{constexpr} vs \texttt{consteval}.
  
  \textbf{C2 (Understand):} Jelaskan mengapa \texttt{std::string\_view} tidak allocation saat substring.
  
  \textbf{C3 (Apply):} Tulis fungsi \texttt{parse\_voltage(std::string\_view)} yang ekstrak nilai double dari \texttt{"V=230.5"}.
  
  \textbf{C4 (Analyze):} Bandingkan \texttt{std::format} vs \texttt{std::ostringstream} untuk performa \& safety.
  
  \textbf{C5 (Evaluate):} Kapan sebaiknya pakai \texttt{consteval} vs \texttt{constexpr} untuk perhitungan jaringan?
  
  \textbf{C6 (Create):} Buat \texttt{consteval} function untuk generate \texttt{std::array} impedansi saluran standar IEC 60909.
\end{frame}

\begin{frame}[allowframebreaks]{Problem Set: Fase 2 (C1--C6)}
  \frametitle{Problem Set: Fase 2 — Memory Management}
  
  \textbf{C1:} Apa itu Rule of Five? Kapan Rule of Zero cukup?
  
  \textbf{C2:} Jelaskan ownership transfer dalam move constructor \texttt{MatrixBuffer}.
  
  \textbf{C3:} Implementasikan \texttt{MatrixBuffer::operator=(MatrixBuffer\&\&)} dengan \texttt{std::exchange}.
  
  \textbf{C4:} Analisis: Mengapa \texttt{get\_row()} return \texttt{std::span} bukan \texttt{T*}?
  
  \textbf{C5:} Evaluasi: Kapan \texttt{unique\_ptr\textlessT[]\textgreater} lebih baik dari \texttt{vector\textlessT\textgreater} untuk matrix?
  
  \textbf{C6:} Desain \texttt{MatrixBuffer} yang support \texttt{operator()(row,col)} dengan bounds check.
\end{frame}

\begin{frame}[allowframebreaks]{Problem Set: Fase 3 (C1--C6)}
  \frametitle{Problem Set: Fase 3 — Concepts \& Generic Programming}
  
  \textbf{C1:} Definisi \texttt{concept} vs \texttt{template} klasik.
  
  \textbf{C2:} Jelaskan \texttt{requires std::invocable\textlessFunc, const EventType\&\textgreater}.
  
  \textbf{C3:} Buat concept \texttt{Serializable} yang require \texttt{to\_json()} method.
  
  \textbf{C4:} Analisis type erasure di \texttt{EventBus::HandlerBase}. Apa trade-off?
  
  \textbf{C5:} Evaluasi: Kapan pakai \texttt{if constexpr} vs template specialization?
  
  \textbf{C6:} Extend \texttt{EventBus} dengan \texttt{unsubscribe} dan priority queue handler.
\end{frame}

\begin{frame}[allowframebreaks]{Problem Set: Fase 4 (C1--C6)}
  \frametitle{Problem Set: Fase 4 — Ranges \& Pipelines}
  
  \textbf{C1:} Apa perbedaan \texttt{std::views::filter} vs \texttt{std::ranges::filter}?
  
  \textbf{C2:} Jelaskan lazy evaluation pada pipeline \texttt{filter | transform | take}.
  
  \textbf{C3:} Tulis pipeline cari top-3 overload tertinggi dari \texttt{vector\textlessBranch\textgreater}.
  
  \textbf{C4:} Analisis: Mengapa \texttt{take(2)} setelah \texttt{transform} hemat performa?
  
  \textbf{C5:} Evaluasi: Kapan \texttt{std::optional} lebih baik dari \texttt{std::expected}?
  
  \textbf{C6:} Buat \texttt{views::chunk_by} custom untuk grup readings per substation.
\end{frame}

\begin{frame}[allowframebreaks]{Problem Set: Fase 5 (C1--C6)}
  \frametitle{Problem Set: Fase 5 — Concurrency}
  
  \textbf{C1:} Perbedaan \texttt{std::thread} vs \texttt{std::jthread}.
  
  \textbf{C2:} Jelaskan \texttt{stop\_token} cooperative cancellation pattern.
  
  \textbf{C3:} Implementasi thread-safe queue dengan \texttt{mutex} + \texttt{condition\_variable}.
  
  \textbf{C4:} Analisis: Kapan \texttt{memory\_order\_relaxed} aman vs butuh \texttt{seq\_cst}?
  
  \textbf{C5:} Evaluasi: \texttt{scoped\_lock} vs \texttt{lock\_guard} untuk deadlock prevention.
  
  \textbf{C6:} Desain work-stealing thread pool untuk N-1 contingency parallellization.
\end{frame}

\begin{frame}[allowframebreaks]{Problem Set: Fase 6 (C1--C6)}
  \frametitle{Problem Set: Fase 6 — Advanced \& Hardware Sympathy}
  
  \textbf{C1:} Keuntungan \texttt{std::expected} vs exception untuk error handling.
  
  \textbf{C2:} Jelaskan monadic operations: \texttt{and\_then}, \texttt{transform}, \texttt{or\_else}.
  
  \textbf{C3:} Refactor \texttt{solve\_power\_flow} pakai \texttt{and\_then} untuk chain validasi.
  
  \textbf{C4:} Analisis cache miss: Hitung L1 miss rate AoS vs SoA untuk stride-64 access.
  
  \textbf{C5:} Evaluasi: Kapan SoA TIDAK lebih cepat dari AoS? (Hint: single-pass streaming)
  
  \textbf{C6:} Desain \texttt{GridSOA} dengan SIMD-friendly alignment (\texttt{alignas(64)}).
\end{frame}

\begin{frame}[allowframebreaks]{Problem Set: Capstone (C1--C6)}
  \frametitle{Problem Set: Capstone — N-1 Contingency Engine}
  
  \textbf{C1:} Jelaskan alur N-1 contingency screening secara ringkas.
  
  \textbf{C2:} Mengapa \texttt{std::span} dipakai untuk split workload ke thread pool?
  
  \textbf{C3:} Modifikasi \texttt{worker\_thread} untuk handle exception di dalam thread.
  
  \textbf{C4:} Analisis race condition pada \texttt{global\_results} tanpa mutex.
  
  \textbf{C5:} Evaluasi: Kapan thread pool size > CPU cores bermanfaat?
  
  \textbf{C6:} Extend engine: Tambah \texttt{N-2} contingency dengan combinatorial generation
  dan early pruning via \texttt{std::ranges}.
\end{frame}

\end{document}